\section{Limitations}
\label{sec:limitations}

This section consolidates the paper's scoping statements rather than distributing hedging throughout the text.

\textbf{Rule-based policies only.}
All experiments use rule-based policies rather than learned policies. The results characterize what happens in the current controller family under the present evaluation protocol. They do not directly generalize to learned controllers, which may exhibit different failure modes and different interactions with memory.

\textbf{Vector observations only.}
The observation space is restricted to vector observations with a deliberately constructed PO mechanism in which the goal is visible at reset and then only flashed periodically. This is a useful controlled setting for isolating the role of memory, but it is not equivalent to realistic visual partial observability.

\textbf{Small task suite.}
The positive conclusions are carried by four solvable tasks (\task{reach-v3}, \task{push-v3}, \task{pick-place-v3}, \task{basketball-v3}), while six contact-rich tasks remain at $0\%$ for all policies. The paper's strongest claims are therefore scoped to the regime of geometrically solvable manipulation under intermittent goal observability.

\textbf{High variance on push-v3.}
\task{push-v3} shows unusually high across-seed variance in both PO and FO evaluations (std $\approx$ 25--29 percentage points). This suggests sensitivity to initialization and execution noise and motivates reporting distributional summaries (std and bootstrap CI) rather than only mean success.

\textbf{Single PO protocol.}
The PO protocol uses a fixed flash mechanism (goal visible at reset, then flashed every $N$ steps for 1 step). Other PO constructions---such as noisy observations, partial state masking, or adversarial observation scheduling---may produce different bottleneck structures.

\textbf{Door-open boundary.}
The \task{door-open-v3} failure is a boundary marker, not a problem that the current approach is close to solving. The perception--execution gap (97.7\% pull entry, 0\% success) indicates that the required fix is at the contact-dynamics level, not at the memory or recovery level. This boundary should be read as a demarcation of scope rather than as a preview of future success.

\textbf{Retry conclusion is local.}
The ablation shows that a naive rollback-heavy retry mechanism can be harmful under the current 200-step horizon on \task{pick-place-v3}. It does not establish that all recovery mechanisms are ineffective, only that recovery cost must be modeled explicitly in finite-horizon settings.
